\section{Aggregation and the owner of the deviation}\label{sec:ownership}
The homogeneous specialization of Section~\ref{sec:heterogeneous} provides
a separate way to distinguish a planning-value gap from a participant's
incentive. Replicate the nominal services to $n$
simultaneous copies of A and $n$ copies of B. Retain the per-bus physical
parameters and fee 7. Provide $n$ early connectors of power 10 kW and $n$
terminal connectors of power 30 kW. Supply scales with the market:
\begin{equation}\label{eq:scaledF}
 F_n(E,L)=4E+\frac{E^2+L^2}{10n}.
\end{equation}
This scaling keeps the marginal-price range bounded; it is an assumption
of the family, not an asymptotic claim for an unchanged supply curve.

If $m$ buses each cover an A--B pair, the remaining services use $n-m$
A-only and $n-m$ B-only buses. Operating cost is $14n-7m$, early energy
is in $[10m,10n]$, and terminal energy is $30n-E$. All terminal charges
fit: each paired bus takes one connector, while each unpaired A/B pair
shares one connector serially. Thus the projected opportunities used below
have explicit physical realizations.
All A-to-B crossmatchings are admitted in this aggregate model. Identical
energies make their projected outcomes depend only on $m$; fixed pair duties
are not being imposed here. In the threshold notation, $r_m=10m$.

\begin{theorem}[Planning gap and whole-owner regret]\label{thm:replication}
The hull value is $\CH_n=7591n/80$. Every physical planner optimum takes
an integer $m$ nearest to $27n/40$, retaining both choices at a tie.
Writing $\delta=m-27n/40$,
\begin{equation}\label{eq:scaledgap}
 \Delta_n=\frac{20\delta^2}{n}\le\frac5n.
\end{equation}
For one owner controlling all services and buses, the own-price regret is
\begin{equation}\label{eq:scaledregret}
 r_n=
 \begin{cases}
 m\,40\delta/n,&\delta\ge0,\\
 (n-m)(-40\delta/n),&\delta<0.
 \end{cases}
\end{equation}
In particular, along $n=40k+1$,
\[
 \Delta_n=\frac{169}{80n}\longrightarrow0,\qquad
 r_n=\frac{351}{40}+\frac{169}{40n}
 \longrightarrow\frac{351}{40}.
\]
For multiples of 40, both quantities are zero.
\end{theorem}

Approximate convexity of an aggregate does not by itself bound a composite
owner's absolute incentive to replace its entire aggregate. This is consistent
with aggregation perspectives \citep{starr1969,hreinsson2021}. In this
family regret per bus and regret relative to physical cost vanish; the
nonvanishing statement concerns the whole owner's \emph{absolute} fixed-price
gain along a declared subsequence.

Now divide ownership among $n$ operators, each responsible for one A/B
pair, with a reserved early connector and a reserved terminal connector.
Each may use one or two buses and choose its two-bus early charge in
$[0,10]$. Those reserved rights ensure unilateral deviations remain feasible
with every other participant's plan held fixed.
An aggregate optimum may pair A and B copies across these ownership labels.
Before making an individual comparison, implement its same projected outcome
by assigning $m$ owners a one-bus plan and the remaining $n-m$ owners a
two-bus plan on their own reserved service pairs. Identical service copies
and the explicit connector assignment make this possible with unchanged
aggregate load, operating cost and marginal price. This is a reimplementation
of the economic outcome, not a claim that a cross-matched literal schedule
already respects the newly imposed ownership rights.
This step is specific to identical services. In the heterogeneous example
of Section~\ref{sec:heterogeneous}, fixed pairs with energies $(5,14)$
and $(15,15)$ require 10 early kWh when both use one bus, whereas the
crossed complete assignment requires only 9. The latter outcome cannot
be transferred to that fixed-pair institution by relabeling. No individual
regret bound for heterogeneous fixed ownership is asserted below.

\begin{corollary}[Reserved participants]\label{cor:individual}
For any physical planner optimum reimplemented in the reserved-pair
institution as above, each reserved pair's regret at its marginal price
is at most $20/n$. The regrets sum to
\eqref{eq:scaledregret}. Along $n=40k+1$, each dissatisfied one-bus owner
can gain only $13/n$.
\end{corollary}
The proof is the same per-pair bill comparison
(Appendix~\ref{app:replication}). At $n=20$ the two tied physical planner
optima with $m=13$ and $m=14$ have whole-owner regrets 7 and 14,
respectively. Neither the whole-owner calculation
nor the reserved-pair calculation specifies an institution for unreserved
shared chargers or strategic anticipation of a tariff reset.

\subsection{Which uniform controls make a small gap informative?}
The preceding example does not imply that planning error and incentives are
unrelated. The selected price matters in the classical link between lost
opportunity cost and duality gaps \citep{huaBaldick2017}. The following
elementary convex-analysis bound makes the relevant uniform controls explicit;
it is an explanatory consequence, not a new general pricing principle.

\begin{proposition}[Gap and own-price regret under uniform curvature]
\label{prop:regretbound}
Let the projected complete-plan set $S=\{(u(x),c(x)):x\in\X\}$ be nonempty
and compact, and let $K$ be its load projection, of diameter $R_K$.
Suppose $F$ is differentiable, $\mu$-strongly convex and has
$\Lambda$-Lipschitz gradient on a relative neighborhood of
$\operatorname{conv}K$, where $0<\mu\le\Lambda$.
Use one fixed Euclidean norm on its affine hull.
For any physical planner optimum $x$, at $p=\nabla F(u(x))$,
\[
 \Delta\le r(x;p)\le
 \Delta+\Lambda R_K\sqrt{2\Delta/\mu}.
\]
If loads lie in a proper affine subspace, gradients and norms can be taken
tangentially: normal price components add the same amount to every bill.
Economic price admissibility is a separate qualification.
\end{proposition}
\begin{proof}
Let $(u^*,c^*)$ minimize $c+F(u)$ on $\operatorname{conv}S$ and put
$q=\nabla F(u^*)$. Hull optimality and linear minimization give
$V(q)=c^*+q\cdot u^*$. With the Bregman difference
$B_F(u,u^*)=F(u)-F(u^*)-q\cdot(u-u^*)$, this yields the exact identity
\[
 c(x)+q\cdot u(x)-V(q)=\Delta-B_F(u(x),u^*)\ge0.
\]
Consequently $\|u(x)-u^*\|\le\sqrt{2\Delta/\mu}$.
For this fixed plan, regret is the maximum of affine price functions with
slopes $u(x)-u(y)$, so it is $R_K$-Lipschitz in price. Transfer from $q$ to
$p$, use gradient smoothness and discard the nonpositive Bregman term to
obtain the upper bound. The lower bound is the tangent inequality from
Proposition~\ref{prop:support}.
\end{proof}

Uniformly bounded $R_K,\Lambda$ and a uniform positive lower bound on
$\mu$ therefore make planner regret vanish as $\Delta\to0$. In our
replication family, however,
$\mu_n=\Lambda_n=1/(5n)$ and $R_{K,n}=10\sqrt2\,n$ in two-load
coordinates. The displayed upper becomes $\Delta_n+40|\delta|$.
Indeed the fixed-$q$ score gap is zero there, so retaining its Bregman
term sharpens this to $40|\delta|$. The exact nonvanishing whole-owner
regret is thus consistent with the bound. Arbitrary nonoptimal plans require
their own excess $c(x)+F(u(x))-\CH$ in place of $\Delta$; nonsmooth supply
or vanishing curvature needs a different argument.
